\documentclass[aspectratio=1610,14pt]{beamer}
\usepackage[ngerman]{babel}
\usepackage{fontspec}
\usepackage{tikz}
\usepackage{calc}
\usetikzlibrary{positioning, arrows.meta}

% ------------------------------------------------------------------
% Typography
% ------------------------------------------------------------------
\setsansfont{Helvetica Neue}[
  UprightFont = * Light,
  BoldFont    = * Bold,
  ItalicFont  = * Light Italic,
  BoldItalicFont = * Bold Italic,
]
\setmonofont{Menlo}[Scale=0.88]
\renewcommand{\familydefault}{\sfdefault}

% ------------------------------------------------------------------
% Palette: slate ink, warm paper, copper accent
% ------------------------------------------------------------------
\definecolor{paper}{HTML}{F7F5F0}
\definecolor{ink}{HTML}{15202B}
\definecolor{mute}{HTML}{6B6E76}
\definecolor{rule}{HTML}{D9D4C7}
\definecolor{copper}{HTML}{B85C2A}
\definecolor{steel}{HTML}{2C4A5C}

% ------------------------------------------------------------------
% Strip beamer chrome
% ------------------------------------------------------------------
\setbeamertemplate{navigation symbols}{}
\setbeamertemplate{footline}{}
\setbeamertemplate{headline}{}
\setbeamertemplate{itemize item}{\color{copper}\textbullet}
\setbeamertemplate{itemize subitem}{\color{copper}\textendash}
\setbeamercolor{normal text}{fg=ink, bg=paper}
\setbeamercolor{structure}{fg=copper}
\setbeamercolor{alerted text}{fg=copper}
\setbeamercolor{itemize item}{fg=copper}
\setbeamerfont{itemize/enumerate body}{size=\normalsize}

% ------------------------------------------------------------------
% Frame title: large, with a thin copper rule and a small slide counter
% ------------------------------------------------------------------
\setbeamerfont{frametitle}{family=\sffamily, series=\bfseries, size=\Large}
\setbeamercolor{frametitle}{fg=ink, bg=paper}
\setbeamertemplate{frametitle}{%
  \vspace*{1.4em}%
  \begin{tikzpicture}[remember picture, overlay]
    \node[anchor=north west, inner sep=0pt, text depth=0pt]
      at ([xshift=0.9cm, yshift=-0.9cm]current page.north west)
      {\usebeamerfont{frametitle}\usebeamercolor[fg]{frametitle}\insertframetitle};
    \draw[copper, line width=1.2pt]
      ([xshift=0.9cm, yshift=-1.95cm]current page.north west)
      -- ++(1.2cm, 0);
    \node[anchor=north east, inner sep=0pt, font=\sffamily\scriptsize, text=mute]
      at ([xshift=-0.9cm, yshift=-1.0cm]current page.north east)
      {\insertframenumber};
  \end{tikzpicture}%
  \vspace{0.4em}%
}

% ------------------------------------------------------------------
% Section page: full-bleed slate with big numeral + section name
% ------------------------------------------------------------------
\AtBeginSection[]{%
  \begingroup
  \setbeamercolor{background canvas}{bg=ink}
  \begin{frame}[plain, noframenumbering]
    \begin{tikzpicture}[remember picture, overlay]
      \node[anchor=west, font=\sffamily\bfseries\fontsize{320}{320}\selectfont,
            text=copper, opacity=0.18]
        at ([xshift=-0.5cm]current page.west)
        {\insertsectionnumber};
      \node[anchor=center, text=paper, font=\sffamily\bfseries\fontsize{32}{38}\selectfont,
            text width=0.78\paperwidth, align=center]
        at (current page.center)
        {\hyphenpenalty=10000\exhyphenpenalty=10000\hyphenchar\font=-1\relax
         \tolerance=9999\emergencystretch=3em\relax\insertsection};
      \node[anchor=south west, font=\sffamily\scriptsize]
        at ([xshift=0.9cm, yshift=0.7cm]current page.south west)
        {\textcolor{copper}{\rule[0.6ex]{1.2cm}{0.8pt}}\textcolor{paper}{\quad Rack'n'Stack 2026 \quad \textbullet \quad Veit Heller}};
    \end{tikzpicture}
  \end{frame}
  \endgroup
}

% ------------------------------------------------------------------
% Title page
% ------------------------------------------------------------------
\setbeamertemplate{title page}{%
  \begin{tikzpicture}[remember picture, overlay]
    % Copper accent block, top-left
    \fill[copper] ([xshift=0cm, yshift=0cm]current page.north west)
      rectangle ([xshift=0.45cm, yshift=-\paperheight]current page.north west);
    % Title
    \node[anchor=north west, text=ink, font=\sffamily\bfseries\fontsize{34}{40}\selectfont,
          text width=0.85\paperwidth, inner sep=0pt]
      at ([xshift=1.4cm, yshift=-1.5cm]current page.north west)
      {\inserttitle};
    % Subtitle
    \node[anchor=north west, text=steel, font=\sffamily\fontsize{16}{20}\selectfont,
          text width=0.78\paperwidth, inner sep=0pt]
      at ([xshift=1.4cm, yshift=-5.0cm]current page.north west)
      {\insertsubtitle};
    % Copper rule between subtitle and credits
    \draw[copper, line width=1.2pt]
      ([xshift=1.4cm, yshift=2.0cm]current page.south west)
      -- ([xshift=1.4cm+1.5cm, yshift=2.0cm]current page.south west);
    % Author block (bottom-left)
    \node[anchor=south west, text=ink, font=\sffamily\bfseries\normalsize, inner sep=0pt]
      at ([xshift=1.4cm, yshift=1.25cm]current page.south west)
      {\insertauthor};
    \node[anchor=south west, text=mute, font=\sffamily\footnotesize, inner sep=0pt]
      at ([xshift=1.4cm, yshift=0.7cm]current page.south west)
      {Port Zero};
    % Conference + date, bottom-right
    \node[anchor=south east, text=mute, font=\sffamily\footnotesize, inner sep=0pt]
      at ([xshift=-1.4cm, yshift=1.25cm]current page.south east)
      {Rack'n'Stack 2026};
    \node[anchor=south east, text=mute, font=\sffamily\footnotesize, inner sep=0pt]
      at ([xshift=-1.4cm, yshift=0.7cm]current page.south east)
      {Nürnberg \textbullet{} 21.\,April 2026};
  \end{tikzpicture}%
}

% Background colour for normal frames
\setbeamercolor{background canvas}{bg=paper}

% ------------------------------------------------------------------
% Slide helpers
% ------------------------------------------------------------------
\newcommand{\hero}[1]{%
  \vfill
  \begin{center}
    {\sffamily #1\par}
  \end{center}
  \vfill
}
\newcommand{\subhero}[1]{%
  \begin{center}
    {\color{mute}\sffamily #1\par}
  \end{center}
}
\newcommand{\copperaccent}[1]{\textcolor{copper}{\textbf{#1}}}
\newcommand{\Statt}{\textcolor{mute}{\itshape Statt:}\,}
\newcommand{\Besser}{\textcolor{copper}{\bfseries Besser:}\,}

\title{Von der Cloud\\ins Rack}
\subtitle{Wie Netzautomatisierung souveräne RZ-Umgebungen ermöglicht}
\date{21.\,April 2026}
\author{Veit Heller}
\institute{Port Zero}
\begin{document}
  \begin{frame}[plain, noframenumbering]
    \titlepage
  \end{frame}

  % ------------------------------------------------------------------
  \section{Das Problem}
  % ~ 4 min
  % ------------------------------------------------------------------
  \begin{frame}{Warum überhaupt zurück?}
    \begin{itemize}
      \item Data Act
      \item NIS2
      \item DORA
      \item AI Act
    \end{itemize}
  \end{frame}

  \begin{frame}{Der Komfort der Hyperscaler}
    \begin{itemize}
      \item Self-Service-Portale
      \item APIs überall
      \item Reproduzierbare Deployments
      \item Klare Verantwortlichkeiten
    \end{itemize}
  \end{frame}

  \begin{frame}{Die Angst}
    \vfill
    \begin{center}
      {\sffamily\itshape „Könnt ihr mir mal ein VLAN machen?“\par}
    \end{center}
    \vfill
    \subhero{Bild im Kopf: zurück in die 2000er.}
  \end{frame}

  \begin{frame}{Die eigentliche Frage}
    \vfill
    \begin{center}
      {\color{mute}\itshape Nicht:}\\[0.2em]
      \textbf{Wie automatisiere ich alles?}

      \vspace{1.4em}
      {\color{copper}\itshape Sondern:}\\[0.2em]
      \textbf{Warum scheitern die meisten Projekte,\\
      obwohl die Technik längst gelöst ist?}
    \end{center}
    \vfill
  \end{frame}

  % ------------------------------------------------------------------
  \section{Die Technik ist das einfache Problem}
  % ~ 5 min
  % ------------------------------------------------------------------
  \begin{frame}{Die Technik ist das einfache Problem}
    \hero{Automatisierungs-Pipelines\\
    bauen wir seit Jahren.}
  \end{frame}

  \begin{frame}{Architektur-Überblick}
    \vspace{0.8em}
    \begin{center}
      \resizebox{0.98\textwidth}{!}{%
      \begin{tikzpicture}[
        box/.style={
          draw=ink, line width=0.8pt, rounded corners=2pt,
          minimum width=2.6cm, minimum height=1.6cm,
          align=center, font=\sffamily, fill=paper, text=ink
        },
        inbox/.style={
          draw=ink, line width=0.8pt, rounded corners=2pt,
          minimum width=2.6cm, minimum height=0.85cm,
          align=center, font=\sffamily, fill=paper, text=ink
        },
        arrow/.style={-{Stealth[length=2.5mm]}, line width=1pt, copper}
      ]
        \node[box] (sot)  {Source\\of Truth};
        \node[box, right=1.3cm of sot]  (auto) {Automation};
        \node[box, right=1.3cm of auto] (pipe) {Pipeline};

        \node[inbox] (op) at ([xshift=-2.8cm, yshift= 0.85cm]sot.west) {Ticket-System};
        \node[inbox, below=0.6cm of op] (cp) {Kundenportal};
        \node[font=\sffamily, text=ink, below=0.4cm of cp] {…};

        \draw[arrow] (op.east) -- (sot.west);
        \draw[arrow] (cp.east) -- (sot.west);

        \draw[arrow] (sot) -- (auto);
        \draw[arrow] (auto) -- (pipe);
      \end{tikzpicture}%
      }
    \end{center}
  \end{frame}

  % --- Source of Truth ---
  \begin{frame}{Source of Truth: DCIM/IPAM}
    Was modelliert wird:
    \begin{itemize}
      \item Standorte, Racks, Geräte
      \item IP-Bereiche, VLANs
      \item Mandanten, Services
    \end{itemize}
    \vspace{1em}
    {\color{mute}\itshape Beispiele: Netbox, Nautobot.}
  \end{frame}

  % --- Automatisierung ---
  \begin{frame}{Automatisierung als Ausführung}
    \begin{itemize}
      \item Desired State, idempotent
      \item Config Compliance
      \item Drift Detection
      \item Pre- und Post-Checks
    \end{itemize}
    \vspace{1em}
    {\color{mute}\itshape Beispiele: Ansible, Nornir.}
  \end{frame}

  % --- Pipelines ---
  \begin{frame}{Pipelines als Kontrollschicht}
    \vfill
    \begin{center}
      Request \,{\color{copper}$\to$}\, Änderung \,{\color{copper}$\to$}\, Review \,{\color{copper}$\to$}\, Rollout
    \end{center}
    \vfill
    {\color{mute}\itshape Merge Request = Change Request. Linting, Dry-Run, Staging. Rollback durch Revert.}
  \end{frame}

  \begin{frame}{Praxis: WOBCOM}
    \vfill
    \begin{center}
      Mit WOBCOM, Stadtwerke Wolfsburg.\\[0.6em]
      Vom Inventar zum ersten\\
      automatisierten Rollout.\\[0.6em]
      In 6 Monaten.
    \end{center}
    \vfill
    {\color{mute}\itshape „Automate yourself within six months“. DENOG 2019, mit Christian Dieckhoff.\\
    \texttt{media.ccc.de/v/denog11-26-}\\
    \texttt{automate-yourself-within-six-months}}
  \end{frame}

  % ------------------------------------------------------------------
  \section{Warum Projekte trotzdem scheitern}
  % ~ 4 min
  % ------------------------------------------------------------------
  \begin{frame}{Trotzdem scheitern die meisten}
    \hero{Die Pipeline steht.\\
    Der Workflow nicht.}
  \end{frame}

  \begin{frame}{Silo: IT für IT}
    \hero{Die IT automatisiert\\
    für die IT.}
  \end{frame}

  \begin{frame}{Falsche Abstraktionsebene}
    \hero{Wir bauen Tooling.\\
    Die Teams haben Prozesse.}
  \end{frame}

  \begin{frame}{Keine Anbindung ans Geschäft}
    \begin{itemize}
      \item Das Projekt lebt im IT-Bereich.
      \item Niemand fragt die Fachteams.
      \item Schmerzhafte Workflows bleiben schmerzhaft.
    \end{itemize}
  \end{frame}

  % ------------------------------------------------------------------
  \section{Daten folgen Produkten und Prozessen}
  % ~ 8 min
  % ------------------------------------------------------------------
  \begin{frame}{Cloud als Produkt}
    \hero{Cloud-Komfort ist Produktdesign.\\
    Die Technik kam danach.}
  \end{frame}

  \begin{frame}{Die Eigenschaften}
    \begin{itemize}
      \item Definierte Produkte und Services
      \item Self-Service, API-first
      \item Reproduzierbarkeit
      \item Nachvollziehbarkeit
      \item Klare Verantwortlichkeiten
    \end{itemize}
  \end{frame}

  \begin{frame}{Was ist euer Produkt?}
    \vfill
    \begin{center}
      \Statt „VLAN einrichten“

      \vspace{1.4em}
      \Besser „Mandantennetz bereitstellen“
    \end{center}
    \vfill
  \end{frame}

  \begin{frame}{Vor der Technik}
    Drei Fragen, bevor irgendetwas automatisiert wird:
    \begin{itemize}
      \item Was sind unsere Produkte und Services?
      \item Was sind die Prozesse dahinter?
      \item Wo lohnt sich Automatisierung, wo nicht?
    \end{itemize}
  \end{frame}

  \begin{frame}{Vom impliziten zum expliziten Prozess}
    \vfill
    \begin{center}
      Aufschreiben \,{\color{copper}$\to$}\, Vereinfachen \,{\color{copper}$\to$}\, Automatisieren
    \end{center}
    \vfill
    {\color{mute}\itshape Automatisierung ist der letzte Schritt, nicht der erste.}
  \end{frame}

  \begin{frame}{}
    \vfill
    \begin{center}
      \begin{minipage}{0.82\paperwidth}
        \centering\itshape
        „Wenn Sie einen scheiß Prozess digitalisieren,\\
        haben Sie einen scheiß digitalen Prozess!“
      \end{minipage}
    \end{center}
    \vfill
    \begin{flushright}
      {\color{mute}— Thorsten Dirks, ehem. Vorstandsvorsitzender Telefónica Deutschland}
    \end{flushright}
  \end{frame}

  \begin{frame}{Wo Automatisierung Sinn macht}
    Kriterien:
    \begin{itemize}
      \item Wie häufig?
      \item Wie fehleranfällig?
      \item Wie reproduzierbar?
    \end{itemize}
  \end{frame}

  \begin{frame}{Nicht alles. Nicht sofort.}
    \begin{itemize}
      \item Was dem Geschäft weh tut, nicht nur der IT.
      \item Die häufigsten 3 bis 5 Vorgänge zuerst.
      \item Sonderfälle bleiben manuell. Bewusst.
      \item Edge Cases später, wenn der Kern steht.
    \end{itemize}
  \end{frame}

  \begin{frame}{Prozesse definieren Services}
    \hero{Der Servicekatalog ist kein Dokument.\\
    Er ist das, was die Pipelines können.}
  \end{frame}

  % ------------------------------------------------------------------
  \section{Compliance als Nebenprodukt}
  % ~ 4 min
  % ------------------------------------------------------------------
  \begin{frame}{Change Management wird einfacher}
    \vfill
    \begin{center}
      \renewcommand{\arraystretch}{1.5}
      \begin{tabular}{r @{\;\textcolor{copper}{$=$}\;} l}
        Merge Request    & Change Record \\
        Review           & Approval \\
        Pipeline-Log     & Ausführungsnachweis \\
        Source of Truth  & Dokumentation \\
      \end{tabular}
    \end{center}
    \vfill
  \end{frame}

  \begin{frame}{Audits und Zertifizierungen}
    \begin{itemize}
      \item ISO 27001, BSI IT-Grundschutz, SOC 2
      \item Source of Truth, Git, Pipeline. Lückenlose Spur.
      \item Drift Detection als kontinuierlicher Nachweis.
    \end{itemize}
  \end{frame}

  \begin{frame}{Dokumentation als Nebenprodukt}
    \hero{Wenn die Source of Truth stimmt,\\
    ist die Doku aktuell.}
  \end{frame}

  % ------------------------------------------------------------------
  \section{Brownfield-Muster}
  % ~ 4 min
  % ------------------------------------------------------------------
  \begin{frame}{Kein Greenfield nötig}
    \hero{Brownfield ist nicht das Problem.\\
    Brownfield ist der Normalfall.}
  \end{frame}

  \begin{frame}{Schritt 1: Inventar und Modell}
    \begin{itemize}
      \item Bestand erfassen.
      \item DCIM/IPAM befüllen.
      \item Muss nicht perfekt sein.
    \end{itemize}
  \end{frame}

  \begin{frame}{Schritt 2: Ein Prozess, Ende zu Ende}
    \begin{itemize}
      \item Einen Prozess sauber durchautomatisieren.
      \item Tief statt breit.
      \item Quick Win, Vertrauen aufbauen.
    \end{itemize}
  \end{frame}

  \begin{frame}{Schritt 3: Ausweiten}
    \vspace{0.3em}
    \begin{itemize}
      \item Pipeline, Review, nächste Prozesse.
      \item Self-Service für Kund:innen und interne Teams.
      \item Das Modell wächst mit.
    \end{itemize}
  \end{frame}

  % ------------------------------------------------------------------
  \section{Fazit}
  % ~ 3 min
  % ------------------------------------------------------------------
  \begin{frame}{Das schwere Problem}
    \hero{Die Technik ist gelöst.\\
    Die Prozesse müssen folgen.}
  \end{frame}

  \begin{frame}{Rückführung als Upgrade}
    \hero{Souveräne Infrastruktur fühlt sich wie Cloud an,\\
    wenn die IT nicht allein automatisiert.}
  \end{frame}
  {
    \setbeamercolor{background canvas}{bg=ink}
    \begin{frame}[plain, noframenumbering]
      \begin{tikzpicture}[remember picture, overlay]
        \node[anchor=west, text=paper, font=\sffamily\bfseries\fontsize{72}{80}\selectfont]
          at ([xshift=1.4cm]current page.west)
          {Danke.};
        \node[anchor=west, text=copper, font=\sffamily\fontsize{28}{34}\selectfont]
          at ([xshift=1.4cm, yshift=-1.8cm]current page.west)
          {Fragen?};
        \node[anchor=south west, text=mute, font=\sffamily\footnotesize]
          at ([xshift=1.4cm, yshift=1.55cm]current page.south west)
          {Veit Heller \textbullet{} Port Zero};
        \node[anchor=south west, text=mute, font=\sffamily\footnotesize]
          at ([xshift=1.4cm, yshift=0.95cm]current page.south west)
          {\texttt{veit@veitheller.de}};
        \node[anchor=south east, text=mute, font=\sffamily\footnotesize]
          at ([xshift=-1.4cm, yshift=1.55cm]current page.south east)
          {Slides};
        \node[anchor=south east, text=mute, font=\sffamily\footnotesize]
          at ([xshift=-1.4cm, yshift=0.95cm]current page.south east)
          {\texttt{github.com/hellerve/talks}};
      \end{tikzpicture}
    \end{frame}
  }
\end{document}
